\section{Related Work}
\label{sec:related-work}

\paragraph{Reasoning-model evaluation.}
Modern LLM evaluation has moved from single-task accuracy toward broad capability suites and explicit tests of reasoning. MMLU and BIG-bench established the value of evaluating models across heterogeneous subjects and tasks \citep{hendrycks2021mmlu,srivastava2023bigbench}, while GSM8K, chain-of-thought prompting, and self-consistency made multi-step reasoning a central evaluation target \citep{cobbe2021gsm8k,wei2022cot,wang2023selfconsistency}. These benchmarks are essential precursors, but they mainly ask whether a model gets an answer right. Our focus is different: when frontier reasoning models fail in professional contexts, we ask which failure mechanism is visible, whether that mechanism is domain-dependent, and whether stated confidence is usable as a risk signal.

The rise of test-time reasoning systems makes this question urgent. OpenAI's o1 system card describes models trained to spend more time on difficult problems, and DeepSeek-R1 reports reasoning capability induced through reinforcement learning \citep{openai2024o1,guo2025deepseekr1}. Claude 3.7 Sonnet and Gemini 2.5 Pro similarly advertise extended or thinking-oriented behavior \citep{anthropic2025claude37,google2025gemini25}. These system reports motivate evaluation of the model class, but they do not provide an independent taxonomy of professional-domain failures. We therefore treat reasoning-model reports as foundational sources for model selection, not as answers to the reliability question.

\paragraph{Professional-domain benchmarks.}
Several domain-specific benchmarks show that high-stakes language tasks require specialized evaluation. In medicine, MedQA/MedMCQA-style datasets and GPT-4 medical evaluations test clinical knowledge and reasoning, while M-ARC is a recent medical abstraction/reasoning benchmark directly relevant to reasoning-model deployment \citep{jin2021diseaseqa,pal2022medmcqa,nori2023gpt4medicine,kim2025marc}. In law, CUAD grounds contract-review answers in expert annotations, and LegalBench and LexGLUE broaden legal reasoning and legal-language evaluation \citep{hendrycks2021cuad,guha2023legalbench,chalkidis2022lexglue}. In finance, FinQA targets numerical reasoning over financial reports, while recent risk-aware and long-context hallucination benchmarks argue that standard financial-agent evaluations can miss deployment-critical errors \citep{chen2021finqa,financebench2025risk,phantom2025finance}. Journalism and fact-checking benchmarks such as ClaimBuster and FEVER focus on check-worthy claims and evidence-backed verification \citep{hassan2017claimbuster,thorne2018fever}. Nuclear-energy evaluation is newer, with emerging work on frontier LLMs for nuclear research and proposed nuclear-engineering benchmarks \citep{almeldein2025nuclear,nucbench2026}.

These studies motivate our five-domain design, but they are not substitutes for it. A medical-only benchmark can reveal clinical failure, and a finance-only benchmark can reveal SEC-filing hallucination, but neither can test whether the dominant error mode changes when the same model family is moved from medicine to law, nuclear review, financial audit, and journalism. Paper 03 holds the coding scheme fixed across domains so that domain-specific failure distributions become measurable rather than anecdotal.

\paragraph{Hallucination, faithfulness, and context failure.}
Hallucination and faithfulness research provides the vocabulary for our feature-hallucination code. Early summarization work showed that fluent generated text can be unfaithful to source documents \citep{maynez2020faithfulness}; later surveys organized hallucination definitions, causes, detection, and mitigation methods for NLG and LLMs \citep{ji2023survey,huang2023hallucinationSurvey}. Recent hallucination benchmarks such as HalluLens further refine the measurement problem \citep{halluLens2025}. Our taxonomy incorporates this literature but does not reduce all failures to hallucination. Professional reasoning failures also include applying an inappropriate learned template, compressing away critical context from long documents, and expressing unjustified confidence.

\paragraph{Confidence calibration.}
Confidence is often proposed as a practical safeguard: if a model knows when it is wrong, a workflow can defer uncertain cases to humans. Prior calibration work complicates that hope. \citet{lin2022calibration} study whether models can express uncertainty in words, \citet{kadavath2022language} argue that language models often contain information about answer correctness, and confidence-elicitation studies show that uncertainty estimates can depend strongly on prompting and measurement format \citep{xiong2023calibration}. We build on this line by measuring confidence miscalibration as a failure mode rather than as a separate post-hoc diagnostic, and by asking whether calibration error remains high across all professional domains.

\paragraph{Interpreting chain-of-thought traces.}
Our method inspects generated reasoning traces, but prior work warns against over-interpreting them. Chain-of-thought prompting and self-consistency can improve performance \citep{wei2022cot,wang2023selfconsistency}, yet generated explanations may be unfaithful to the process that produced the answer \citep{turpin2023unfaithful}. The recent argument to stop anthropomorphizing intermediate tokens reinforces this caution \citep{stopanthro2024}. Accordingly, we use traces as observable artifacts for coding failure symptoms, not as privileged access to a model's internal cognition. This conservative stance lets the taxonomy remain behavioral and reproducible: annotators label what the model asserted, omitted, overgeneralized, or overconfidently endorsed.

\paragraph{Positioning.}
The closest related works each cover one axis of our study: broad reasoning benchmarks, single-domain professional benchmarks, hallucination taxonomies, confidence calibration, or model-card safety reports. None, as located in this literature pass, jointly evaluates frontier reasoning models across medicine, law, nuclear safety review, financial audit, and journalism using a common failure taxonomy and calibration analysis. Our contribution is therefore a comparative reliability map: the same models, the same F1--F4 codes, and the same annotation protocol applied across domains where model errors have materially different consequences.
